\chapter{Apêndice C — Estrutura de Pastas e Reprodutibilidade}
\label{apendiceC}

Este apêndice descreve a organização lógica dos arquivos utilizados na
dissertação, bem como os princípios adotados para garantir reprodutibilidade
empírica, transparência metodológica e rastreabilidade dos resultados
apresentados nos capítulos principais.

\section{Estrutura de diretórios}

Os arquivos do projeto foram organizados de forma modular, separando dados,
scripts, tabelas e figuras. A estrutura geral segue o padrão abaixo:

\begin{verbatim}
project_root/
│
├── data/
│   ├── raw/                # Dados brutos (preços, IV, DVOL, etc.)
│   ├── processed/          # Dados tratados e alinhados
│
├── notebooks/
│   ├── dvol_pipeline.ipynb # Construção da IV e BVRP
│   ├── rv_calculation.ipynb
│   ├── ml_models.ipynb
│
├── scripts/
│   ├── utils_vol.py        # Funções auxiliares (RV, defasagens)
│   ├── ml_pipeline.py
│
├── tables/
│   ├── tab_ols_*.tex
│   ├── tab7/
│   ├── tab8/
│
├── figs/
│   ├── cap6/
│   ├── cap7/
│   ├── cap8/
│
├── chapters/
│   ├── cap1_introducao.tex
│   ├── cap2_referencial.tex
│   ├── ...
│
└── main.tex
\end{verbatim}

Essa organização permite separar claramente:
(i) dados brutos e tratados;
(ii) rotinas computacionais;
(iii) resultados finais utilizados no documento.

\section{Reprodutibilidade empírica}

Todos os resultados empíricos apresentados na dissertação são gerados a partir
de scripts determinísticos, com controle explícito de sementes aleatórias
(\textit{random seeds}) nos modelos de aprendizado de máquina, quando aplicável.
As figuras e tabelas reportadas no texto principal são exportadas diretamente
dos \textit{notebooks} e scripts Python para arquivos \texttt{.png} e
\texttt{.tex}, evitando intervenções manuais.

As estimações seguem rigorosamente a ordenação temporal das informações,
assegurando ausência de \textit{look-ahead bias} e consistência fora da amostra,
conforme discutido nos Capítulos~\ref{chap:metodologia},
\ref{chap:cap6_ml} e \ref{chap:ml_bvrp}.

\section{Disponibilidade do código}

O código completo utilizado na dissertação encontra-se organizado e documentado
pelo autor. Por razões de licenciamento de dados e uso de fontes proprietárias
(plataformas de derivativos de criptoativos), os dados brutos e parte dos scripts
não são disponibilizados publicamente. No entanto, todas as etapas metodológicas,
definições de variáveis e especificações empíricas são descritas de forma
detalhada no texto e nos apêndices, permitindo replicação conceitual e
implementação independente por terceiros.

\section{Resumo}

Este apêndice documenta a estrutura computacional subjacente à dissertação,
reforçando o compromisso com boas práticas de pesquisa empírica, organização de
projetos quantitativos e reprodutibilidade dos resultados.
